\section{Phương trình hóa học và tính theo phương trình}

Sau khi đổi dữ kiện về các đại lượng thích hợp, mối liên hệ giữa các chất được xác định bởi \textbf{tỉ lệ hệ số trong phương trình hóa học đã cân bằng}.

\begin{ghinho}
Quy trình cơ bản:
\begin{enumerate}
\item Viết và cân bằng phương trình hóa học.
\item Đổi dữ kiện về số mol nếu cần.
\item Dùng tỉ lệ hệ số để tìm số mol chất cần tính.
\item Đổi số mol về đại lượng đề yêu cầu.
\item Kiểm tra đơn vị và độ lớn kết quả.
\end{enumerate}
\end{ghinho}

\begin{vidu}
Cho \SI{13}{\gram} Zn phản ứng hoàn toàn với dung dịch \ce{HCl} dư:
\[
\ce{Zn + 2HCl -> ZnCl2 + H2}.
\]
Tính thể tích \ce{H2} thu được ở điều kiện \(V_m=\SI{24.8}{\liter\per\mole}\). Cho \(M_{\ce{Zn}}=65\) g/mol.
\end{vidu}
\begin{loigiai}
\[
n_{\ce{Zn}}=\frac{13}{65}=\SI{0.20}{\mole}.
\]
Theo phương trình, \(n_{\ce{H2}}=\SI{0.20}{\mole}\), nên
\[
V_{\ce{H2}}=0.20\times24.8=\SI{4.96}{\liter}.
\]
\end{loigiai}

\begin{vidu}
Đốt cháy hoàn toàn \SI{1.5}{\ton} quặng sphalerite chứa \(80\%\) \ce{ZnS}:
\[
\ce{2ZnS + 3O2 -> 2ZnO + 2SO2}.
\]
Tính khối lượng \ce{ZnO} tối đa thu được. Cho \(M_{\ce{ZnS}}=97\), \(M_{\ce{ZnO}}=81\) g/mol.
\source{Phỏng theo bài luyện Zn từ quặng sphalerite, đề chuyên Hóa Lâm Đồng 2026.}
\end{vidu}
\begin{loigiai}
\[
m_{\ce{ZnS}}=1.5\times0.80=\SI{1.20}{\ton}.
\]
Vì tỉ lệ mol \(\ce{ZnS}:\ce{ZnO}=1:1\):
\[
m_{\ce{ZnO}}=1.20\times\frac{81}{97}\approx\SI{1.00}{\ton}.
\]
\end{loigiai}

\begin{casiobox}
Trong bài nhiều bước, nên giữ một biểu thức tính hoàn chỉnh rồi mới bấm máy để giảm sai số do làm tròn trung gian. Ví dụ:
\[
1.5\times0.80\times\frac{81}{97}.
\]
Nếu cần thử nhiều giá trị phần trăm quặng khác nhau, có thể nhập biểu thức theo biến rồi dùng \key{CALC}.
\end{casiobox}

\begin{vidu}
Đốt cháy hoàn toàn \SI{5.4}{\gram} Al:
\[
\ce{4Al + 3O2 -> 2Al2O3}.
\]
Tính khối lượng \ce{Al2O3} thu được. Cho \(M_{\ce{Al}}=27\), \(M_{\ce{Al2O3}}=102\) g/mol.
\end{vidu}
\begin{loigiai}
\[
n_{\ce{Al}}=\frac{5.4}{27}=\SI{0.20}{\mole},\qquad
n_{\ce{Al2O3}}=\frac{2}{4}\times0.20=\SI{0.10}{\mole}.
\]
\[
m_{\ce{Al2O3}}=0.10\times102=\SI{10.2}{\gram}.
\]
\end{loigiai}

\begin{tuluyen}
Cho \SI{11.2}{\gram} Fe phản ứng hoàn toàn với \ce{HCl} dư. Tính số mol và thể tích \ce{H2} ở \(V_m=\SI{24.8}{\liter\per\mole}\). Cho \(M_{\ce{Fe}}=56\) g/mol.
\dapso{\SI{0.20}{\mole}; \SI{4.96}{\liter}.}
\end{tuluyen}

\begin{tuluyen}
Nung \SI{500}{\kilo\gram} \ce{CaCO3} tinh khiết: \(\ce{CaCO3 -> CaO + CO2}\). Tính khối lượng \ce{CaO} tối đa. Cho \(M_{\ce{CaCO3}}=100\), \(M_{\ce{CaO}}=56\) g/mol.
\dapso{\SI{280}{\kilo\gram}.}
\end{tuluyen}

\begin{tuluyen}
Từ \SI{2.00}{\ton} \ce{Fe2O3} tinh khiết, khử hoàn toàn bằng CO: \(\ce{Fe2O3 + 3CO -> 2Fe + 3CO2}\). Tính khối lượng Fe tối đa. Cho \(M_{\ce{Fe2O3}}=160\), \(M_{\ce{Fe}}=56\) g/mol.
\dapso{\SI{1.40}{\ton}.}
\end{tuluyen}

\begin{tuluyen}
Dùng \SI{800}{\kilo\gram} \ce{ZnS} tinh khiết để sản xuất Zn theo hai bước \ce{ZnS -> ZnO -> Zn}. Giả sử các phản ứng hoàn toàn. Tính khối lượng Zn. Cho \(M_{\ce{ZnS}}=97\), \(M_{\ce{Zn}}=65\) g/mol.
\dapso{\(\approx\SI{536}{\kilo\gram}\).}
\end{tuluyen}
